% RESEARCH_DRAFT_SHA256: 96f62da4001703bffa875021715d455b46ca9878eeaa95553ce28ec7492d6845
% Layout and source-reference wrappers expanded for the web renderer.
\documentclass[11pt]{article}
\usepackage[margin=25mm]{geometry}
\usepackage{fontspec,unicode-math}
\setmainfont{Latin Modern Roman}
\setmathfont{Latin Modern Math}
\usepackage{amsmath,amssymb,mathtools}
\usepackage{xurl,hyperref}
\hypersetup{colorlinks=true,urlcolor=blue,linkcolor=blue}
\setcounter{secnumdepth}{0}
\setlength{\parindent}{0pt}
\setlength{\parskip}{5pt}
\setlength{\emergencystretch}{3em}
\newcommand{\sref}[2]{\hyperlink{source:#1:#2}{[S#2]}}
\newcommand{\cataloguescope}[1]{\paragraph{Source question.}#1}
\newcommand{\floor}[1]{\left\lfloor#1\right\rfloor}
\begin{document}
\section*{Research-attempt provenance}
This standalone research attempt retains the complete original definition body below,
including its historical source and preparation notes. Those notes describe the
earlier catalogue preparation; the new attempt and its verification are identified
separately after that body. No original ID, catalogue status, or source statement is changed.

% UnsolvedMath ID 30003380; source catalogue code OWR-15210-006
\section{Average-Case Hardness of Quasi-Cyclic Code Problems}\label{30003380}
{\small UnsolvedMath ID 30003380 · Computer Science}
\subsection{Definitions and mathematical statement}

\paragraph{Shared notation.}Standard number systems are denoted by $\mathbb N,\mathbb Z,\mathbb Q,\mathbb R,\mathbb C$; index ranges, computational models and specialized notation follow the cited problem source and the definitions below.


An element of $R_p$ is represented by a binary polynomial modulo $x^p-1$. Its Hamming weight counts nonzero coefficients; total weight sums the weights of the two displayed polynomial components.

A code $C\subseteq\Sigma^N$ is a set of permitted length-$N$ words. Its Hamming minimum distance is the minimum number of unequal coordinates of two distinct words. For a linear $[N,K,D]_q$ code, rate is $K/N$ and relative distance is $D/N$; decoding recovers a codeword from corrupted observations.
\paragraph{Statement recorded in the supplied source.}
Let $R_p=\mathbb F_2[x]/(x^p-1)$. Extend average-case decoding-hardness results to quasi-cyclic codes, and prove a polynomial equivalence between the following problems: finding $e_0,e_1\in R_p$ of bounded total Hamming weight with $e_0+e_1h=s$, and deciding whether there exist nonzero $h_0,h_1\in R_p$ of bounded total Hamming weight with $h_0+h_1h=0$. \href{https://doi.org/10.4171/owr/2017/5}{[S2]}
\paragraph{Scope qualification.}
The distribution of $h,s$, weight bounds, promise, and reduction direction are source-specific; polynomial equivalence is asked rather than assumed.
\paragraph{Source status.}
The supplied catalogue labels this record partially solved. Its source-status metadata was checked against the supplied snapshot on 2026-10-11. This import does not independently establish current open status or mathematical correctness. The catalogue literature assessment is dated 2026-08-21.
\subsection{Short English statement}
Extend decoding hardness to quasi-cyclic codes and relate the two stated sparse polynomial problems.
\subsection{Sources}
\begin{itemize}
\item[S1] Ulam AI and the UnsolvedMath Contributors, UnsolvedMath: A Curated Collection of Open Mathematics Problems, record 30003380 (OWR-15210-006). Catalogue curation and metadata: CC BY 4.0. The supplied source question is retained; definitions and the English summary are editorial additions.
\url{https://huggingface.co/datasets/ulamai/UnsolvedMath}.
\item[S2] Cryptography (2018), Oberwolfach Reports, DOI 10.4171/owr/2017/5, pp. 213--266
\url{https://doi.org/10.4171/owr/2017/5}.
\item[S3] On the Hardness of Decoding Quasi-Cyclic Codes
\url{https://www.scitepress.org/PublishedPapers/2026/142964/}.
\item[S4] Efficient Encryption from Random Quasi-Cyclic Codes
\url{https://arxiv.org/abs/1612.05572}.
\textit{Underlying Oberwolfach Reports material retains its source-specific licence, including CC BY-SA 4.0 where applicable; see THIRD\_PARTY\_NOTICES.md.}
\end{itemize}
\paragraph{Preparation and review.}
Original question attributed to the cited source and UnsolvedMath contributors. Codex (GPT-6; exact serving revision not exposed) assisted with definitions, English summaries, source inspection and formatting. No new human mathematical authorship or independent proof verification is asserted. Canonical ID 30003380 and source code OWR-15210-006 checked against the complete supplied registry and existing repository statements; known flagged repetitions and obvious identical conjectures excluded. Full mathematical-equivalence/literature review has not been performed. Statement preparation and final review followed contribution prompts C and D; automated checks verify structure and identity, not open status.
\label{end:30003380}

\section{Literature and context check}
The QC-MDPC talk in OWR 5/2017, pp.~241--243, was read, including its exact Problems 2 and 3. Problem 3 requires both displayed components nonzero, while Problem 2 permits zero error components. The source asks for average-case hardness and polynomial equivalence; it does not supply the distributions needed for a distribution-preserving reduction. The following algebra and explicit five-bit counterexample test a natural homogenization route, not those full claims.
\section{Research attempt: Coset decoding and a failed quasi-cyclic homogenization}

\paragraph{The two geometric tasks.}
In $R_p=\mathbb F_2[x]/(x^p-1)$ define
\[
 C_h=\{(uh,u):u\in R_p\}\subseteq\mathbb F_2^{2p}.
\]
This is a binary linear code of dimension $p$, since the second coordinate recovers $u$. It is invariant under simultaneous cyclic shifts because multiplication by $x$ commutes with multiplication by $h$. Every solution of $e_0+e_1h=s$ is uniquely
\[
 (e_0,e_1)=(s+uh,u),\qquad u=e_1.
\]
Thus decoding asks whether the affine coset $(s,0)+C_h$ contains a vector of weight at most $t$, whereas the codeword question concerns low-weight vectors of $C_h$ with both coordinates nonzero. If $h$ is a unit, any nonzero $u$ has $uh\ne0$, so the latter is the ordinary nonzero minimum-distance question. For general $h$ that extra condition matters: at $h=0$, Problem 3 is always negative, while decoding reduces to testing $|s|\le t$, by choosing $u=0$ and noting that $|s|+|u|$ cannot be smaller than $|s|$.

A direct exact decoding algorithm enumerates $u$ of weight at most $t$, computes $s+uh$, and checks the total weight. This is complete because any feasible solution has $|u|\le t$. Naive cyclic multiplication costs $O(p^2)$ bit operations, giving
\[
 O\left(p^2\sum_{j=0}^{\min(t,p)}\binom pj\right)
\]
up to comparable weight-checking work. For fixed $t$ this is polynomial; for growing weights it gives no uniform polynomial algorithm. The codeword problem has the analogous enumeration with the two nonzero checks. This does not establish average-case hardness or an equivalence.

\paragraph{Ordinary homogenization already needs a promise.}
Adjoin $(s,0)$ to the binary linear span of $C_h$, obtaining
\[
 C'=C_h\cup\bigl((s,0)+C_h\bigr)
\]
when $(s,0)\notin C_h$. A short nonzero vector in $C'$ might already lie in $C_h$, unrelated to the desired syndrome. Only a promise that the old code has minimum distance greater than $t$ removes this false positive. Moreover the dimension usually becomes $p+1$, outside the original systematic dimension-$p$ family. Closing under cyclic shifts creates a further obstruction, even when the old-distance promise holds.

\paragraph{An explicit failure after cyclic closure.}
Take $p=5$, $h=1$, $s=1+x+x^2$, and $t=1$. The old code is $C_1=\{(u,u)\}$, whose nonzero vectors have weight at least two. The decoding instance is negative because
\[
 |e_0|+|e_1|\ge |e_0+e_1|=|s|=3.
\]
Here the inequality holds coefficientwise: a nonzero coefficient of the sum requires at least one nonzero summand coefficient.

But in $R_5$ one has
\[
 (x+x^2+x^4)(1+x+x^2)=1.
\]
Expansion in $\mathbb F_2[x]$ gives $x+x^5+x^6$, since each of $x^2,x^3,x^4$ occurs twice. Reduction by $x^5=1$ and $x^6=x$ gives $1$. Therefore the smallest cyclic submodule containing $(s,0)$ contains $(1,0)$, of weight one. Adding it to $C_1$ and using $(1,1)\in C_1$ also gives $(0,1)$, so the cyclic closure is actually all of $R_5^2$. A procedure that homogenizes by adjoining all cyclic shifts would turn this negative decoding instance into a code with trivial low-weight words and would leave the original code family. The weight-one vectors $(1,0)$ and $(0,1)$ do not satisfy Problem 3's requirement that both components be nonzero. Thus the example is a false positive for ordinary minimum-distance homogenization and a departure from the original family, not a positive instance of original Problem 3 at weight one.

This counterexample defeats that particular reduction. It does not show that no more elaborate polynomial reduction between the source problems exists; such a reduction would need to preserve the family, its nonzero-component convention, the weight thresholds and the intended average-case distributions.
\section{Verification and adversarial audit}
The code dimension, affine-coset parametrization, unit versus nonunit boundary and exhaustive-search completeness were checked directly. The five-bit inverse was computed by enumerating all $32$ ring elements with exact cyclic binary multiplication; the unique inverse had exponent set $\{1,2,4\}$. The no-decoding-at-weight-one conclusion follows independently from the weight inequality. The distinction between an arbitrary nonzero codeword and both components nonzero is retained.
\paragraph{Audit corrections.} Kept both-component nonzero requirements explicit and did not substitute a generic minimum-distance problem without a unit hypothesis. Added an exact cyclic-closure counterexample rather than asserting that a general linear-code reduction automatically remains quasi-cyclic.
\paragraph{External computation.} Python 3.14.2 enumerated the $32$ elements of $R_5$ by bit masks and multiplied modulo cyclic index five. For $s=7$ the unique inverse mask was $22$, corresponding to $x+x^2+x^4$; the standalone script \texttt{check\_cyclic\_inverse.py} also verified coset minimum weight three and old-code minimum weight two. No cryptographic parameter experiment or hardness test was run.
\section{Scope and remaining gaps}
The original average-case hardness statement and polynomial equivalence remain unproved. The attempted homogenization fails before distributional issues even arise: it admits old-code words, can change dimension, and its cyclic closure can introduce spurious very short vectors. A family- and distribution-preserving reduction with controlled weights is the first missing step.
\paragraph{Integration review.} Codex (GPT-6; exact serving revision not exposed) reviewed the written arguments and expanded layout and source-reference wrappers. Repository metadata and the production web renderer were checked using Python 3.13 and JavaScript with MathJax 3.2.2. This model-assisted review is not independent human proof verification; no PDF compilation or formal proof checking is asserted.

\section{Completion Estimate}
\noindent\textbf{COMPLETION: 0\%}
The calculations below the preserved source statement investigate a concrete route,
but establish no substantiated new progress toward the entire original target.
Known special cases and failed sufficient conditions are not credited as a solution.
\section{References}
The original source references and catalogue attribution are preserved in the
Sources section above. The following URLs were actually consulted for this attempt;
retention of an additional original reference does not claim fresh consultation of it.
\begin{itemize}
\item \url{https://ems.press/content/serial-article-files/46669} (consulted 2026-10-11). 
\end{itemize}
\section{Model and contributor disclosure}
OpenAI gpt-6-astra, requested reasoning configuration ultra, generated this research
attempt on 2026-10-11. The exact serving revision is not exposed. Codex GPT-6
performed orchestration, repository-format review, and written-proof review of this attempt; its exact serving revision
is also not exposed. The mathematical route, boundary checks and final adversarial
self-audit followed contribution prompts A then B in
\texttt{docs/contribution\_prompts.md}, after reading \texttt{CONTRIBUTING.md}.
Python 3.14.2 assembled the standalone documents; any mathematical computation is
specified in Verification. No theorem prover was used. No human mathematical
authorship, substantial human proof correction, or independent human proof checking
is asserted. Model-assisted review, source retrieval, and structural checks do not
establish correctness or novelty. Full private reasoning transcripts are unavailable
as reproducibility artifacts; the complete arguments, inputs and outcomes actually
used are stated here or in the companion batch audit.
\section{Final status}
Unresolved as a new research contribution. This is a concrete route investigation,
not a claim of new complete resolution or a current-status certification. A known
answer reproduced for context does not become a new contribution. The catalogue's
source-status label is preserved; this attempt status does not assert that a
source-labelled solved question is open.
\end{document}
